\section{Introduction}
\label{sec:intro}
Many outputs of machine-learning systems are compositions: vectors of non-negative parts that sum to one.
Softmax class probabilities, topic proportions and species mixes are such vectors, as the quantity to be predicted or as the prediction itself.
Conformal prediction \citep{vovk2005algorithmic,angelopoulos2023gentle} turns any predictor into a set or region with finite-sample marginal coverage under exchangeability, and it is widely used for classifier outputs \citep{romano2020classification,angelopoulos2021uncertainty}.
For a response that itself lies on the simplex, a practitioner must choose a nonconformity score, and the natural candidates come from different geometries: Euclidean and $L_1$ distances, the Hellinger distance, divergences, the Aitchison distance of log-ratio analysis \citep{aitchison1982statistical}, and the likelihood under a Dirichlet model.

Marginal coverage is guaranteed for all of them, so it cannot decide among them.
What distinguishes scores is the size of the region and how evenly the coverage is spread over covariates and over the simplex, in particular near its boundary.
Exact conditional coverage is impossible in a distribution-free setting \citep{vovk2012conditional,barber2021limits}, so papers report relaxations that are \emph{estimated} from a finite test set: worst-slab coverage \citep{cauchois2020knowing}, size-stratified coverage \citep{angelopoulos2021uncertainty}, class-conditional gaps \citep{ding2023class}.
How reliable these estimates are is the question we ask, in settings where the true conditional coverage and the oracle region are known.

\paragraph{Research questions.}
\textbf{RQ1} Does the coverage of split conformal prediction follow its exact finite-sample law for every score on the simplex, and what does this imply for the sizes needed to certify coverage?
\textbf{RQ2} At equal marginal coverage, how do score geometries differ in region size and in conditional and boundary coverage, and under which noise models is each preferable?
\textbf{RQ3} How reliable are the estimators of worst-cell coverage used to compare scores: how large is their bias, does it depend on the method, and can it reverse a ranking?
\textbf{RQ4} Under which conditions do the findings stop holding (covariate shift, rounding and ties, a misspecified mean, the calibration scheme)?
\textbf{RQ5} Do the conclusions carry over to classifier outputs and their set-valued predictors?

\paragraph{Contributions.}
To our knowledge, and within the sources we consulted (Section~\ref{sec:related}), the following are new in combination.
\begin{itemize}\setlength\itemsep{0pt}
\item A known-truth benchmark for conformal regions on the simplex: 15 scores from Euclidean to Dirichlet-likelihood, two noise families, 20 configurations, with ground-truth conditional coverage and volume ratios to the oracle highest-density region (Section~\ref{sec:exp_geom}).
\item An exact expression for the expected value of the worst-cell coverage estimate (Proposition~\ref{prop:min}), checked against simulation, and evidence on how the bias differs between methods, how often it reverses a pair of scores as a function of their true gap, and how many replications make a ranking reproducible (Section~\ref{sec:exp_est}).
\item Negative results reported as found: a half-split and a plug-in correction of the worst-cell estimator do not improve ranking, and a Mahalanobis score gives no gain over the Aitchison score.
\item A transfer to classifier sets (LAC, APS, RAPS, Mondrian) including a CIFAR-10 CNN, an evaluation protocol, and conditions (covariate shift, rounding, a misspecified mean, CV+/jackknife+) where results change (Sections~\ref{sec:exp_clf}--\ref{sec:discussion}).
\end{itemize}
The exact coverage law (Proposition~\ref{prop:beta}) is a known property of split conformal prediction; we use it as the reference.

\paragraph{Scope.}
Data are synthetic with a known truth, except one CIFAR-10 network; we do not claim that rankings carry over to real compositional data.
